\documentclass[10pt,a4paper]{amsart}
\usepackage{amsmath,amssymb,amsthm,graphicx}
\usepackage{verbatim}
\newenvironment{reference}{\comment}{\endcomment}
%\newenvironment{reference}{}{}
\newenvironment{slogan}{\comment}{\endcomment}
\newenvironment{history}{\comment}{\endcomment}

% For commutative diagrams we use Xy-pic
\usepackage[all]{xy}

% We use 2cell for 2-commutative diagrams.
\xyoption{2cell}
% Theorem environments.
%
\theoremstyle{plain}
\newtheorem{theorem}[subsection]{Theorem}
\newtheorem{proposition}[subsection]{Proposition}
\newtheorem{lemma}[subsection]{Lemma}

\theoremstyle{definition}
\newtheorem{definition}[subsection]{Definition}
\newtheorem{example}[subsection]{Example}
\newtheorem{exercise}[subsection]{Exercise}
\newtheorem{situation}[subsection]{Situation}

\theoremstyle{remark}
\newtheorem{remark}[subsection]{Remark}
\newtheorem{remarks}[subsection]{Remarks}

\numberwithin{equation}{subsection}

% Macros
%
\def\lim{\mathop{\mathrm{lim}}\nolimits}
\def\colim{\mathop{\mathrm{colim}}\nolimits}
\def\Spec{\mathop{\mathrm{Spec}}}
\def\Hom{\mathop{\mathrm{Hom}}\nolimits}
\def\Ext{\mathop{\mathrm{Ext}}\nolimits}
\def\SheafHom{\mathop{\mathcal{H}\!\mathit{om}}\nolimits}
\def\SheafExt{\mathop{\mathcal{E}\!\mathit{xt}}\nolimits}
\def\Sch{\mathit{Sch}}
\def\Mor{\mathop{\mathrm{Mor}}\nolimits}
\def\Ob{\mathop{\mathrm{Ob}}\nolimits}
\def\Sh{\mathop{\mathit{Sh}}\nolimits}
\def\NL{\mathop{N\!L}\nolimits}
\def\CH{\mathop{\mathrm{CH}}\nolimits}
\def\proetale{{pro\text{-}\acute{e}tale}}
\def\etale{{\acute{e}tale}}
\def\QCoh{\mathit{QCoh}}
\def\Ker{\mathop{\mathrm{Ker}}}
\def\Im{\mathop{\mathrm{Im}}}
\def\Coker{\mathop{\mathrm{Coker}}}
\def\Coim{\mathop{\mathrm{Coim}}}

% Boxtimes
%
\DeclareMathSymbol{\boxtimes}{\mathbin}{AMSa}{"02}

%
% Macros for moduli stacks/spaces
%
\def\QCohstack{\mathcal{QC}\!\mathit{oh}}
\def\Cohstack{\mathcal{C}\!\mathit{oh}}
\def\Spacesstack{\mathcal{S}\!\mathit{paces}}
\def\Quotfunctor{\mathrm{Quot}}
\def\Hilbfunctor{\mathrm{Hilb}}
\def\Curvesstack{\mathcal{C}\!\mathit{urves}}
\def\Polarizedstack{\mathcal{P}\!\mathit{olarized}}
\def\Complexesstack{\mathcal{C}\!\mathit{omplexes}}
% \Pic is the operator that assigns to X its picard group, usage \Pic(X)
% \Picardstack_{X/B} denotes the Picard stack of X over B
% \Picardfunctor_{X/B} denotes the Picard functor of X over B
\def\Pic{\mathop{\mathrm{Pic}}\nolimits}
\def\Picardstack{\mathcal{P}\!\mathit{ic}}
\def\Picardfunctor{\mathrm{Pic}}
\def\Deformationcategory{\mathcal{D}\!\mathit{ef}}
\usepackage{cleveref}
\crefname{originalreference}{original reference}{original references}
\Crefname{originalreference}{Original reference}{Original references}
\makeatletter
\newlabel{more-algebra-lemma-homotopy-koszul-abstract}{{0626}{source}{}{sourcepointer-1}{}}
\newlabel{more-algebra-lemma-homotopy-koszul-abstract@cref}{{[originalreference][0][]0626}{[1][0][]source}{}{}{}}
\makeatother
\title{Poincare lemma for divided power polynomial rings and their p-adic completions}
\author{The Stacks Project Authors}
\date{stacks-project:crystalline.tex:local-snapshot-20261012; GNU FDL 1.2 or later}
\begin{document}
\maketitle
\noindent\textit{Editorial note.} All delivered mathematical text is the author's own, copied verbatim from the pinned Stacks Project chapter file crystalline.tex; no statement or proof step was rewritten, reordered or removed.\par\medskip
\noindent\textit{Editorial note.} Span recomputation by label pairing (LANE\_BRIEF 10.2): in this chapter the author's proof environment follows the statement environment as a standalone block; the spec covers the complete environment pair (statement span, proof span) as recorded in the pool candidate.\par\medskip
\ifdefined\hypertarget\hypertarget{sourcepointer-1}{}\fi
\noindent\textit{Source reference.} Original reference 0626: sources/stacks-chow/tags.txt, source lines 7785--7785. Cross-chapter standard prerequisite cited by the author inside the counted span; it is used as a pointer and its own proof is not part of this item.\par\medskip
\begin{lemma}
\label{lemma-poincare}
Let $A$ be a ring. Let $P = A\langle x_i \rangle$ be a divided
power polynomial ring over $A$. For any $A$-module $M$ the complex
$$
0 \to M \to
M \otimes_A P \to
M \otimes_A \Omega^1_{P/A, \delta} \to
M \otimes_A \Omega^2_{P/A, \delta} \to \ldots
$$
is exact. Let $D$ be the $p$-adic completion of $P$.
Let $\Omega^i_D$ be the $p$-adic completion of the $i$th exterior
power of $\Omega_{D/A, \delta}$. For any $p$-adically complete
$A$-module $M$ the complex
$$
0 \to M \to
M \otimes^\wedge_A D \to
M \otimes^\wedge_A \Omega^1_D \to
M \otimes^\wedge_A \Omega^2_D \to \ldots
$$
is exact.
\end{lemma}

\begin{proof}
It suffices to show that the complex
$$
E :
(0 \to A \to P \to \Omega^1_{P/A, \delta} \to
\Omega^2_{P/A, \delta} \to \ldots)
$$
is homotopy equivalent to zero as a complex of $A$-modules.
For every multi-index $K = (k_i)$ we can consider the subcomplex $E(K)$
which in degree $j$ consists of
$$
\bigoplus\nolimits_{I = \{i_1, \ldots, i_j\} \subset \text{Supp}(K)}
A
\prod\nolimits_{i \not \in I} x_i^{[k_i]}
\prod\nolimits_{i \in I} x_i^{[k_i - 1]}
\text{d}x_{i_1} \wedge \ldots \wedge \text{d}x_{i_j}
$$
Since $E = \bigoplus E(K)$ we see that it suffices to prove each of the
complexes $E(K)$ is homotopic to zero. If $K = 0$, then
$E(K) : (A \to A)$ is homotopic to zero. If $K$ has nonempty (finite)
support $S$, then the complex $E(K)$ is isomorphic to the complex
$$
0 \to A \to
\bigoplus\nolimits_{s \in S} A \to
\wedge^2(\bigoplus\nolimits_{s \in S} A) \to
\ldots \to \wedge^{\# S}(\bigoplus\nolimits_{s \in S} A) \to 0
$$
which is homotopic to zero, for example by
More on Algebra, Lemma \ref{more-algebra-lemma-homotopy-koszul-abstract}.
\end{proof}

\end{document}
